% thm-fluctuations -- anchor CERW.Frozen.fluctuation_bounds. Source: cerw.tex:186-213 (label thm:fluctuations); proof fluctuations.tex:1-302 (flat 845-1146).
\paragraph{Paper (verbatim).} ``For every $p>0$, there are constants $C,n_0>0$, depending only on
$d,\eps,p$, such that for every integer $n\geq n_0$, with probability at least $1-Cn^{-p}$,
$B(0,(a-CQ)N)\cap\Z^d\subseteq A_n\subseteq V_n\subseteq B(0,(a+CQ^{1/d}L)N)\cap\Z^d$, where
$L=\log(n+2)$. On the same event, $|(D_n/N)\triangle B(0,a)|+|\frac{|A_n|}{N^d}-\omega_da^d|\leq CQ$ and
$\sup_{x\in\Z^d}|\frac{\ell_n(x)}N-2d\eps(a-|x|/N)_+|\leq CQ^{1/d}$. All these bounds hold eventually
almost surely, with deterministic constants and a finite random starting time.''
\paragraph{Proof.} Fix $p$ and intersect the events of \texttt{prop-coarse}, \texttt{lem-local},
\texttt{lem-radial}, \texttt{s-localmart}, \texttt{s-vector} and \texttt{s-quadratic} (the stopped
quadratic martingale). The complement has measure at most $Cn^{-p}$, and the rest is deterministic.
Let $W=NQ^{1/d}$.
\emph{Setup (\texttt{s-global}).} Choose a deterministic $K>a+1$ with $D_n\subseteq B(0,(K-1)N)$
(\texttt{prop-coarse}, \texttt{s-cell-norm}). eq:approx with $M_n\le CN$ gives
$|\widetilde\ell_n-U_{D_n}|\le\delta_0$ on $B(0,KN)$, with $\delta_0=C\sqrt NL$ ($d=2$) or
$C(\sqrt{NL}+L)$.
\emph{Contact (\texttt{s-contact-cell}, \texttt{s-contact-bound}).} Let $b$ be the inradius of $D_n$,
$E=D_n\setminus B(0,b)$ and $m=|E|$. There are $y_0$ with $|y_0|=b$ and an unoccupied lattice site
$z$ with $|z-y_0|\le\sqrt d/2$ and $|z|\ge b$. Since $U_{B(0,b)}(y_0)=0$ (\texttt{lem-geometry}) and
$u_v\cdot(v-y_0)|v-y_0|^{-d}\ge2^{1-d}|v|^{1-d}$ for $|v|\ge b$, we get $cm/N^{d-1}\le U_{D_n}(y_0)$.
By \texttt{s-cellmodulus}, \texttt{s-pointwise}, \texttt{s-localmart} and $\ell_n(z)=0$,
$U_{D_n}(y_0)\le C(\sqrt{B_zL}+L)$. By \texttt{lem-geometry} (sup bound for $E$), $|U_E|\le Cm^{1/d}$.
\emph{Variance ($d\ge3$; \texttt{s-var-high}).} At occupied sites,
$\ell_n\le c_0s_b+C(m^{1/d}+L)+C\sqrt{B_yL}$, with $s_b=(b-|\cdot|)_+$ and $B_y=(k*\ell_n)(y)$. Young
turns this into $\ell_n\le c_0s_b+C(V+L)+\lambda k*\ell_n$ everywhere. With
$k*s_b\le k_0s_b+J$ (\texttt{Generic.Lattice.sum\_mul\_le\_of\_lipschitz}) and the moments
$k_0\le C$, $J\le CL$ (\texttt{s-weight-sums}), the resolvent bound
(\texttt{Generic.Lattice.le\_of\_le\_add\_conv}) gives $\ell_n\le C(s_b+V+L)$. Then $B_z\le C(V+L)$,
and the contact inequality with \texttt{Generic.Young.le\_of\_le\_mul\_rpow\_mul\_sqrt\_add} gives
$m\le CN^dQ$.
\emph{Variance ($d=2$; \texttt{s-var-planar}).} First $m\le CN^{3/2}L$. Then
$\ell_n\le c_0(b-|x|)_++CN^{3/4}L^{1/2}+C\sqrt NL$, so $B_z\le CN$ (\texttt{s-weight-sums}), and
$m\le CN^{3/2}\sqrt L=CN^2Q$. Hence $\ell_n\le c_0(b-|y|)_++CW$.
\emph{Shell (\texttt{s-shellW}).} $M_{\rm sh}(r)\le CW$ for $r\ge b+3$.
\emph{Radius (\texttt{s-quadratic}, \texttt{s-inradius}).} $|X_n|^2=n-2\eps\sum_{A_n}|x|+\mathcal Q_n$ and
$|\mathcal Q_n|\le C(N\sqrt{nL}+NL)$. With $|\sum_{A_n}|x|-\int_D|v||\le CN^d$,
$\int_D|v|=\frac{d\omega_d}{d+1}b^{d+1}+O(Nm)$ and $|X_n|\le CN$, we get
$|1-\frac{c_0\omega_d}{d+1}(b/N)^{d+1}|\le CQ$. Since $c_0\omega_da^{d+1}/(d+1)=1$
(\texttt{Guards.limitRadius\_spec}), \texttt{Generic.Young.abs\_sub\_le\_of\_abs\_one\_sub\_mul\_pow}
gives $|b/N-a|\le CQ$.
\emph{Conclusions (\texttt{s-volume-profile}).} $D=B(0,b)\sqcup E$ gives the volume bound and the
inner inclusion (after enlarging $C$). On $B(0,KN)$,
$|\widetilde\ell_n-c_0(aN-|y|)_+|\le C(\delta_0+NQ+NQ^{1/d})\le CW$, and outside both sides vanish.
This is the profile bound.
\emph{Outer (\texttt{s-tailend}, \texttt{s-outer}).} $F(b)\le CNQ$. Halving with $K(A)=C(W+1)$
(\texttt{Generic.Halving}) gives $F(b+K_0WL)\le\Lambda=CN^{2-d}L$. If the path reached radius
$b+3K_1WL$, the last-entrance interval above $\beta=b+2h$ would have $m_{\rm cap}\le CNL$ sites.
\texttt{s-crossing} then gives $(h-1)^2\le C(N^{4/3}L^{8/3}+NL^4)$ for $d=2$, or
$C(N^{2d/(2d-1)}L^{4d/(2d-1)}+NL^3)$ for $d\ge3$, both $o(h^2)$. Hence
$V_n\subseteq B(0,b+3h)\subseteq B(0,(a+CQ^{1/d}L)N)$.
\emph{Eventually a.s. (\texttt{s-assembly}).} With $p=2$, $\sum_nCn^{-2}<\infty$, and
\texttt{ae\_eventually\_notMem} (Borel--Cantelli, no measurability needed) gives the second conjunct
with $C=C(d,\eps,2)$.
\paragraph{Transcription notes.} One declaration, two conjuncts (R-6). $\mathrm{Good}$ is a
\texttt{let} inside the hashed block. Volumes are in $\R_{\ge0}^\infty$, and $\sup$ is written as
$\forall x$ (R-5). The same $C$ serves the event and the probability. $C,n_0$ come after $\eps,p$
and before the realization. $A_n\subseteq V_n$ is kept literally.
